% Source: https://artofproblemsolving.com/community/c3h1352616_medium_probability_problem
% Stable identity: aops-topic-1352616-q01
Let $q$ be the probability of seeing no car in a $5$-minute interval. The $20$ minutes are four disjoint $5$-minute intervals with the same, independent behaviour, so the probability of no car in $20$ minutes is $q^4$. Hence
\[
q^4=1-\dfrac{325}{625}=\dfrac{300}{625}=\dfrac{12}{25},\qquad q=\sqrt[4]{\dfrac{12}{25}}=\dfrac{\sqrt[4]{12}}{\sqrt5}.
\]
The required probability is
\[
1-q=\boxed{1-\sqrt[4]{\dfrac{12}{25}}}\approx0.168 .
\]
(With a nicer number such as $\dfrac{609}{625}$ the same method gives $1-\dfrac{2}{5}=\dfrac{3}{5}$.) The answer is not $\dfrac{1}{4}\cdot\dfrac{325}{625}$: probabilities of ``at least one'' events do not scale linearly with time.
